% !TEX spellcheck = en_US



% ~~~~~~~~~~~~~~~~~~~~~~~~~~~~~~~~~~~~~~~ V
% (NC) 2026 Haluk Bingol
% github.com/halukbingol/LaTeX-Templates
%
% Licensees may copy, distribute, display, and perform the work and make derivative works 
% and remixes based on it only for non-commercial purposes. 
% ~~~~~~~~~~~~~~~~~~~~~~~~~~~~~~~~~~~~~~~ A




% v20261123T2200

\documentclass[aspectratio=169,10pt,compress]{beamer}




% ~~~~~~~~~~~~~~~~~~~~~~~~~~~~~~~~~~~~~~~~~~~~~~~~~~~~~~~~~~~~~~~~~ V hblib
	\usepackage{../../hblib/hblibPresentation} % lib presentation: theme, i/N, \hSection, algorithm2e, ...
	\usepackage{../../hblib/hblibCodeoutput} % lib code + output: \lstcode, \lstterminal, ...
	\usepackage{../../hblib/hblibDatastructures} % lib data structures: \hString, \hStack, \hQueue, \hArray, \hGrid
% ~~~~~~~~~~~~~~~~~~~~~~~~~~~~~~~~~~~~~~~~~~~~~~~~~~~~~~~~~~~~~~~~~ A hblib




% +++++++++++++++++++++++++++++++++++++++ V specific
\DontPrintSemicolon % in algorithms
% +++++++++++++++++++++++++++++++++++++++ A specific



\title[Exceptions, Tests: Junior]{
	Exception Handling and Unit Testing\\
	Junior
}
\subtitle{
	\cisOOO\ -- Level 3: advanced features, performance, larger systems
}
\author[Erkovan]{Sait Erkovan}

\institute[FCIS]{
	Faculty of Computer and Information Sciences\\
	Yeditepe University
}
\date{{\tiny \hbTimeStamp}}

\begin{document}




% ~~~~~~~~~~~~~~~~~~~~~~~~~~~~~~~~~~~~~~ V frm
\begin{frame}[t,fragile]
	\titlepage
\end{frame}
% ~~~~~~~~~~~~~~~~~~~~~~~~~~~~~~~~~~~~~~ A frm




% ~~~~~~~~~~~~~~~~~~~~~~~~~~~~~~~~~~~~~~ V frm
\begin{frame}[t,fragile] \frametitle{Outline}
	\begin{columns}[T]
		\begin{column}{0.48\textwidth}
			\tableofcontents[sections={1-3}, hideallsubsections]
		\end{column}
		\begin{column}{0.48\textwidth}
			\tableofcontents[sections={4-7}, hideallsubsections]
		\end{column}
	\end{columns}

	\vspace{1em}
	\small It builds on \hCode{-2-sophomore}: own exception types, chaining, fixtures and parameterized tests are assumed.
	Interfaces and lambdas are used.
\end{frame}
% ~~~~~~~~~~~~~~~~~~~~~~~~~~~~~~~~~~~~~~ A frm




% ~~~~~~~~~~~~~~~~~~~~~~~~~~~~~~~~~~~~~~ V
\hSection[Cost]{What exceptions cost}
% ~~~~~~~~~~~~~~~~~~~~~~~~~~~~~~~~~~~~~~ A




% ~~~~~~~~~~~~~~~~~~~~~~~~~~~~~~~~~~~~~~ V frm
\begin{frame}[t,fragile] \frametitle{Catch or check? Measured}
	\lstcode[basicstyle=\ttfamily\fontsize{6.4}{7.6}\selectfont, firstline=22, lastline=44, firstnumber=22]{codeoutput/ExceptionCost.java}{Java}
	\only<2->{\begin{textblock*}{0.50\textwidth}(7.6cm,3.2cm)
		\lstoutput[basicstyle=\ttfamily\scriptsize]{codeoutput/ExceptionCost.txt}
		\footnotesize One million words, half of them not numbers; the second of two rounds (the first warms up the JIT).
		The times differ from machine to machine.
	\end{textblock*}}
\end{frame}
% ~~~~~~~~~~~~~~~~~~~~~~~~~~~~~~~~~~~~~~ A frm




% ~~~~~~~~~~~~~~~~~~~~~~~~~~~~~~~~~~~~~~ V frm
\begin{frame}[t,fragile] \frametitle{Where the time goes: the stack trace}
	\lstcode[basicstyle=\ttfamily\fontsize{6.4}{7.6}\selectfont, firstline=24, lastline=30, firstnumber=24]{codeoutput/TraceCost.java}{Java}
	\uncover<2->{\lstoutput[basicstyle=\ttfamily\scriptsize]{codeoutput/TraceCost.txt}}

	\begin{itemize}\footnotesize
		\item<2-> Creating an exception records every frame of the stack: the deeper the call, the slower
		\item<2-> The protected constructor \hCode{(message, cause, enableSuppression, writableStackTrace)} can skip the trace; then the
			  report has no lines to read. Only for very special, measured cases
		\item<2-> ``Without'' still contains the cost of the calls themselves (about 1 ns each)
	\end{itemize}
\end{frame}
% ~~~~~~~~~~~~~~~~~~~~~~~~~~~~~~~~~~~~~~ A frm




% ~~~~~~~~~~~~~~~~~~~~~~~~~~~~~~~~~~~~~~ V frm
\begin{frame}[t,fragile] \frametitle{Rules of thumb}


% +++++++++++++++++++++++++++++++++++++++ V tbl
%: -Tbl. cost rules
\begin{table}
		\footnotesize
		\begin{tabular}{@{}L{0.40\textwidth}L{0.52\textwidth}@{}}
			\toprule
			\textbf{Situation} & \textbf{Prefer} \\
			\midrule
			bad input is common (a file of user data)    & check first (\hCode{hasNextInt}, a pattern), no exception \\
			bad input is rare                            & just parse and catch: simple and fast enough \\
			``not found'' is a normal answer             & a return value (\hCode{-1}, \hCode{Optional}), not an exception \\
			a loop that ends by an exception             & never: end loops with their condition \\
			an exception in a hot loop, measured slow    & check before, or move the check out of the loop \\
			\bottomrule
		\end{tabular}
\end{table}%
% +++++++++++++++++++++++++++++++++++++++ A tbl


	\vspace{0.3em}
	\small Throwing is cheap compared with reading a file or a network; it only matters in tight loops.
	Measure before you optimize.
\end{frame}
% ~~~~~~~~~~~~~~~~~~~~~~~~~~~~~~~~~~~~~~ A frm




% ~~~~~~~~~~~~~~~~~~~~~~~~~~~~~~~~~~~~~~ V
\hSection[Pitfalls]{Pitfalls of finally and resources}
% ~~~~~~~~~~~~~~~~~~~~~~~~~~~~~~~~~~~~~~ A




% ~~~~~~~~~~~~~~~~~~~~~~~~~~~~~~~~~~~~~~ V frm
\begin{frame}[t,fragile] \frametitle{return in finally}
	\lstcoderightstep[0.68][basicstyle=\ttfamily\scriptsize, firstline=7, lastline=21, firstnumber=7][basicstyle=\ttfamily\scriptsize]{codeoutput/FinallyReturn.java}{Java}{codeoutput/FinallyReturn.txt}

	\footnotesize \uncover<2->{A \hCode{return} in \hCode{finally} wins over everything: over the \hCode{return 1}, and over the exception, which
	vanishes without a trace. javac compiles it silently (\hCode{-Xlint:finally} warns); Eclipse warns
	``finally block does not complete normally''. Never return or throw from \hCode{finally}.}
\end{frame}
% ~~~~~~~~~~~~~~~~~~~~~~~~~~~~~~~~~~~~~~ A frm




% ~~~~~~~~~~~~~~~~~~~~~~~~~~~~~~~~~~~~~~ V frm
\begin{frame}[t,fragile] \frametitle{Two problems at once: suppressed exceptions}
	\lstcode[basicstyle=\ttfamily\fontsize{6.4}{7.6}\selectfont, firstline=2, lastline=17, firstnumber=2]{codeoutput/Suppressed.java}{Java}
	\uncover<2->{\lstoutput[basicstyle=\ttfamily\fontsize{6.4}{7.6}\selectfont]{codeoutput/Suppressed.txt}}

	\footnotesize \uncover<2->{The first problem is kept, the second attached to it; a \hCode{finally} with \hCode{close()} would lose the first.}
\end{frame}
% ~~~~~~~~~~~~~~~~~~~~~~~~~~~~~~~~~~~~~~ A frm




% ~~~~~~~~~~~~~~~~~~~~~~~~~~~~~~~~~~~~~~ V frm
\begin{frame}[t,fragile] \frametitle{Checked exceptions and lambdas}
	\lstoutput[basicstyle=\ttfamily\fontsize{6.4}{7.6}\selectfont]{codeoutput/LambdaChecked.txt}
	\footnotesize A lambda for \hCode{forEach} may not throw a checked exception: \hCode{Consumer.accept} declares none. One way out:

	\lstcode[basicstyle=\ttfamily\fontsize{6.4}{7.6}\selectfont, firstline=10, lastline=23, firstnumber=10]{codeoutput/LambdaWrapped.java}{Java}
	\only<2->{\begin{textblock*}{0.48\textwidth}(9.0cm,6.3cm)
		\lstoutput[basicstyle=\ttfamily\fontsize{6}{7.2}\selectfont]{codeoutput/LambdaWrapped.txt}
	\end{textblock*}}
\end{frame}
% ~~~~~~~~~~~~~~~~~~~~~~~~~~~~~~~~~~~~~~ A frm




% ~~~~~~~~~~~~~~~~~~~~~~~~~~~~~~~~~~~~~~ V
\hSection[Errors]{When the JVM is in trouble}
% ~~~~~~~~~~~~~~~~~~~~~~~~~~~~~~~~~~~~~~ A




% ~~~~~~~~~~~~~~~~~~~~~~~~~~~~~~~~~~~~~~ V frm
\begin{frame}[t,fragile] \frametitle{StackOverflowError and OutOfMemoryError}
	\begin{columns}[T]
		\begin{column}{0.54\textwidth}
			\lstcode[basicstyle=\ttfamily\fontsize{6}{7.2}\selectfont, firstline=4, lastline=15, firstnumber=4]{codeoutput/Errors/Deep.java}{Java}
		\end{column}
		\begin{column}{0.44\textwidth}
			\lstcode[basicstyle=\ttfamily\fontsize{6}{7.2}\selectfont, firstline=4, lastline=11, firstnumber=4]{codeoutput/Errors/Hog.java}{Java}
		\end{column}
	\end{columns}
	\uncover<2->{\lstterminal[basicstyle=\ttfamily\fontsize{6}{7.2}\selectfont]{codeoutput/Errors.txt}}
\end{frame}
% ~~~~~~~~~~~~~~~~~~~~~~~~~~~~~~~~~~~~~~ A frm




% ~~~~~~~~~~~~~~~~~~~~~~~~~~~~~~~~~~~~~~ V frm
\begin{frame}[t,fragile] \frametitle{Errors are not for catching}
	\begin{itemize}\small
		\item \hCode{-Xss} sets the stack size, \hCode{-Xmx} the heap; small values make the errors come fast. The number of calls
			  depends on the machine; week 12: a recursion without a base case ends like \hCode{Deep}
		\item An \hCode{Error} means the JVM itself cannot go on as planned: no stack, no memory, a broken class file
		\item \hCode{Deep} catches one only to show it; after a \hCode{StackOverflowError} the program's state may be half-changed
		\item After an \hCode{OutOfMemoryError} even the handler may need memory it does not get
		\item \hCode{catch (Throwable t)} or \hCode{catch (Error e)} belongs, if anywhere, at the very top: log and stop
		\item The fix is in the code or the configuration: a base case, a smaller data structure, a larger \hCode{-Xmx}
	\end{itemize}
\end{frame}
% ~~~~~~~~~~~~~~~~~~~~~~~~~~~~~~~~~~~~~~ A frm




% ~~~~~~~~~~~~~~~~~~~~~~~~~~~~~~~~~~~~~~ V
\hSection[Design]{Exceptions in larger programs}
% ~~~~~~~~~~~~~~~~~~~~~~~~~~~~~~~~~~~~~~ A




% ~~~~~~~~~~~~~~~~~~~~~~~~~~~~~~~~~~~~~~ V frm
\begin{frame}[t,fragile] \frametitle{Layers translate at their borders}
	\begin{center}
		\begin{tikzpicture}[ly/.style={draw, rounded corners=2pt, minimum width=3.3cm, minimum height=0.8cm, align=center,
				font=\small, fill=white}, ex/.style={-{Latex[length=2mm]}, thick, dashed, colorStop},
				nt/.style={font=\scriptsize, anchor=west, align=left}]
			\node[ly] (ui)  at (0,3.0) {user interface};
			\node[ly] (svc) at (0,2.0) {grade service};
			\node[ly] (dao) at (0,1.0) {score store};
			\node[ly, fill=gray!15] (os) at (0,0.0) {files, database};
			\draw[ex] (2.0,0.0) -- (2.0,1.0);
			\draw[ex] (2.0,1.0) -- (2.0,2.0);
			\draw[ex] (2.0,2.0) -- (2.0,3.0);
			\draw[colorStop, thick] (1.65,0.0) -- (2.0,0.0); \draw[colorStop, thick] (1.65,1.0) -- (2.0,1.0);
			\draw[colorStop, thick] (1.65,2.0) -- (2.0,2.0); \draw[colorStop, thick] (1.65,3.0) -- (2.0,3.0);
			\node[nt] at (2.2,0.5) {\hCode{NoSuchFileException}, \hCode{SQLException}};
			\node[nt] at (2.2,1.5) {\hCode{ScoreStoreException} (own type, cause kept)};
			\node[nt] at (2.2,2.5) {the same, or a \hCode{GradeException} with the student's name};
			\node[nt] at (2.2,3.4) {a sentence for the user; the whole chain to the log};
		\end{tikzpicture}
	\end{center}

	\begin{itemize}\footnotesize
		\item Each layer speaks its own language: the interface does not need to know about SQL or file names
		\item Translate at the border of a layer, with the cause (sophomore); handle in one place, at the top
		\item Changing the database then changes only the lowest layer, not its callers
	\end{itemize}
\end{frame}
% ~~~~~~~~~~~~~~~~~~~~~~~~~~~~~~~~~~~~~~ A frm




% ~~~~~~~~~~~~~~~~~~~~~~~~~~~~~~~~~~~~~~ V frm
\begin{frame}[t,fragile] \frametitle{An alternative: Optional}
	\lstcoderightstep[0.66][basicstyle=\ttfamily\fontsize{6.2}{7.3}\selectfont, firstline=5, lastline=23, firstnumber=5][basicstyle=\ttfamily\scriptsize]{codeoutput/OptionalParse.java}{Java}{codeoutput/OptionalParse.txt}

	\footnotesize \uncover<2->{A box that holds a value or nothing; the caller must open it on purpose (\hCode{isPresent}, \hCode{orElse}).
	Good when absence is normal; for real problems, an exception.}
\end{frame}
% ~~~~~~~~~~~~~~~~~~~~~~~~~~~~~~~~~~~~~~ A frm




% ~~~~~~~~~~~~~~~~~~~~~~~~~~~~~~~~~~~~~~ V frm
\begin{frame}[t,fragile] \frametitle{Checked exceptions: the debate}


% +++++++++++++++++++++++++++++++++++++++ V tbl
%: -Tbl. checked exceptions debate
\begin{table}
		\footnotesize
		\begin{tabular}{@{}L{0.46\textwidth}L{0.46\textwidth}@{}}
			\toprule
			\textbf{For} & \textbf{Against} \\
			\midrule
			the signature documents what can go wrong       & a new exception deep down changes every method above it \\
			the compiler does not let you forget a case      & lambdas and streams cannot throw them (previous section) \\
			good for problems that callers can recover from & often caught only to be wrapped or ignored \\
			\bottomrule
		\end{tabular}
\end{table}%
% +++++++++++++++++++++++++++++++++++++++ A tbl


	\vspace{0.4em}
	\begin{itemize}\small
		\item Java is nearly alone: C\#, Kotlin, Scala and Python have no checked exceptions
		\item Java's own newer libraries lean to unchecked (\hCode{UncheckedIOException}, \hCode{java.time} exceptions)
		\item A common team rule: checked for recoverable problems at a module's border, unchecked inside
	\end{itemize}
\end{frame}
% ~~~~~~~~~~~~~~~~~~~~~~~~~~~~~~~~~~~~~~ A frm




% ~~~~~~~~~~~~~~~~~~~~~~~~~~~~~~~~~~~~~~ V
\hSection[Doubles]{Testing code with dependencies}
% ~~~~~~~~~~~~~~~~~~~~~~~~~~~~~~~~~~~~~~ A




% ~~~~~~~~~~~~~~~~~~~~~~~~~~~~~~~~~~~~~~ V frm
\begin{frame}[t,fragile] \frametitle{A class that depends on the outside world}
	\begin{columns}[T]
		\begin{column}{0.42\textwidth}
			\lstcode[basicstyle=\ttfamily\fontsize{6.4}{7.6}\selectfont]{codeoutput/Doubles/ScoreSource.java}{Java}
			\footnotesize \hCode{Report} does not create its source: it gets one (line 7). This is \hDefined{dependency injection}:
			a test can pass in a source it controls.
		\end{column}
		\begin{column}{0.56\textwidth}
			\lstcode[basicstyle=\ttfamily\fontsize{6}{7.2}\selectfont, firstline=7, lastline=27, firstnumber=7]{codeoutput/Doubles/Report.java}{Java}
		\end{column}
	\end{columns}
\end{frame}
% ~~~~~~~~~~~~~~~~~~~~~~~~~~~~~~~~~~~~~~ A frm




% ~~~~~~~~~~~~~~~~~~~~~~~~~~~~~~~~~~~~~~ V frm
\begin{frame}[t,fragile] \frametitle{Stubs, written by hand}
	\lstcode[basicstyle=\ttfamily\fontsize{6.4}{7.6}\selectfont, firstline=9, lastline=26, firstnumber=9]{codeoutput/Doubles/ReportTest.java}{Java}

	\footnotesize Each lambda is a whole \hCode{ScoreSource}: one method, fixed answers. The last one fails on purpose.
\end{frame}
% ~~~~~~~~~~~~~~~~~~~~~~~~~~~~~~~~~~~~~~ A frm




% ~~~~~~~~~~~~~~~~~~~~~~~~~~~~~~~~~~~~~~ V frm
\begin{frame}[t,fragile] \frametitle{A spy, and the results}
	\lstcode[basicstyle=\ttfamily\fontsize{6.4}{7.6}\selectfont, firstline=28, lastline=37, firstnumber=28]{codeoutput/Doubles/ReportTest.java}{Java}
	\uncover<2->{\lstoutput[basicstyle=\ttfamily\scriptsize, firstline=3]{codeoutput/Doubles.txt}}

	\footnotesize A spy also records how it was used: \hCode{Report} must read its source once, not once per number.
\end{frame}
% ~~~~~~~~~~~~~~~~~~~~~~~~~~~~~~~~~~~~~~ A frm




% ~~~~~~~~~~~~~~~~~~~~~~~~~~~~~~~~~~~~~~ V frm
\begin{frame}[t,fragile] \frametitle{Test doubles}


% +++++++++++++++++++++++++++++++++++++++ V tbl
%: -Tbl. test doubles
\begin{table}
		\footnotesize
		\begin{tabular}{@{}L{0.12\textwidth}L{0.42\textwidth}L{0.38\textwidth}@{}}
			\toprule
			\textbf{Kind} & \textbf{Does} & \textbf{In ReportTest} \\
			\midrule
			stub  & returns fixed answers                           & \hCode{fake}, \hCode{broken} \\
			spy   & a stub that also records how it was used        & \hCode{spy} counts the calls \\
			fake  & a simple working version (a list in memory instead of a database) & -- \\
			mock  & checks the calls it expects, usually with a library (Mockito) & -- \\
			\bottomrule
		\end{tabular}
\end{table}%
% +++++++++++++++++++++++++++++++++++++++ A tbl


	\vspace{0.3em}
	\begin{itemize}\small
		\item The failing source tests the \hCode{catch} of \hCode{Report}: the hardest path to reach with a real server
		\item Fast and repeatable: no network, no files, no waiting
		\item With an interface of one method, a lambda is enough; no library needed
	\end{itemize}
\end{frame}
% ~~~~~~~~~~~~~~~~~~~~~~~~~~~~~~~~~~~~~~ A frm




% ~~~~~~~~~~~~~~~~~~~~~~~~~~~~~~~~~~~~~~ V frm
\begin{frame}[t,fragile] \frametitle{Files in tests: @TempDir}
	\lstcode[basicstyle=\ttfamily\fontsize{6}{7.2}\selectfont, firstline=13, lastline=34, firstnumber=13]{codeoutput/Files/ScoreFileTest.java}{Java}
	\only<2->{\begin{textblock*}{0.36\textwidth}(9.8cm,2.4cm)
		\lstoutput[basicstyle=\ttfamily\scriptsize, firstline=3]{codeoutput/Files.txt}
		\footnotesize JUnit makes a new folder for every test and deletes it afterwards. \hCode{badLine} (lines 36--40) writes ``8S''.
	\end{textblock*}}
\end{frame}
% ~~~~~~~~~~~~~~~~~~~~~~~~~~~~~~~~~~~~~~ A frm




% ~~~~~~~~~~~~~~~~~~~~~~~~~~~~~~~~~~~~~~ V
\hSection[Quality]{How good are the tests?}
% ~~~~~~~~~~~~~~~~~~~~~~~~~~~~~~~~~~~~~~ A




% ~~~~~~~~~~~~~~~~~~~~~~~~~~~~~~~~~~~~~~ V frm
\begin{frame}[t,fragile] \frametitle{Coverage}
	\begin{itemize}\small
		\item A coverage tool records which lines and branches the tests ran; Java's usual one is JaCoCo
		\item In Eclipse (EclEmma, built on JaCoCo, if installed): \menu{Coverage As > JUnit Test}; green lines ran, red did not,
			  yellow: a branch only partly taken
		\item With Maven or Gradle, the JaCoCo plugin writes an HTML report after \hCode{mvn test} or \hCode{gradle test}
		\item Useful to find code that \emph{no} test runs (the \hCode{else} of \hCode{Shipping} in the examples deck)
		\item Not a measure of quality: \hCode{WeakTest} on the next frames runs every \hCode{return} of \hCode{letterGrade},
			  and still misses seven of eight small bugs
	\end{itemize}

	\vspace{0.5em}
	\small \hRemark{Not run here}: coverage tools are not part of this deck's sessions.
\end{frame}
% ~~~~~~~~~~~~~~~~~~~~~~~~~~~~~~~~~~~~~~ A frm




% ~~~~~~~~~~~~~~~~~~~~~~~~~~~~~~~~~~~~~~ V frm
\begin{frame}[t,fragile] \frametitle{Mutation testing, by hand}
	\lstcode[basicstyle=\ttfamily\fontsize{6}{7.2}\selectfont, firstline=4, lastline=26, firstnumber=4]{codeoutput/Mutants/mutate.sh}{bash}
\end{frame}
% ~~~~~~~~~~~~~~~~~~~~~~~~~~~~~~~~~~~~~~ A frm




% ~~~~~~~~~~~~~~~~~~~~~~~~~~~~~~~~~~~~~~ V frm
\begin{frame}[t,fragile] \frametitle{What the mutants say}
	\lstterminal[basicstyle=\ttfamily\fontsize{6.4}{7.6}\selectfont]{codeoutput/Mutants.txt}

	\begin{itemize}\footnotesize
		\item A \hDefined{mutant} is a copy of the code with one small change, the kind a programmer makes by mistake
		\item \hCode{WeakTest} (one middle score per letter) notices only one of eight: green tests, weak tests
		\item \hCode{StrongTest} (every boundary and just outside, as in week 11) kills all eight
		\item The tool \hCode{PIT} (pitest.org) does this automatically for Java, with hundreds of mutants
	\end{itemize}
\end{frame}
% ~~~~~~~~~~~~~~~~~~~~~~~~~~~~~~~~~~~~~~ A frm




% ~~~~~~~~~~~~~~~~~~~~~~~~~~~~~~~~~~~~~~ V frm
\begin{frame}[t,fragile] \frametitle{Property-based testing}
	\lstcode[basicstyle=\ttfamily\fontsize{6.4}{7.6}\selectfont, firstline=5, lastline=23, firstnumber=5]{codeoutput/PropertySort.java}{Java}
	\only<2->{\begin{textblock*}{0.46\textwidth}(8.4cm,6.2cm)
		\lstoutput[basicstyle=\ttfamily\scriptsize]{codeoutput/PropertySort.txt}
	\end{textblock*}}

	\footnotesize Instead of a few chosen cases: many random inputs and a \hDefined{property} that must always hold, here ``the same as
	\hCode{Arrays.sort}''. The bug of \hCode{sort} (line 30, \hCode{j > 0}) never moves an element to index 0. Libraries: jqwik.
\end{frame}
% ~~~~~~~~~~~~~~~~~~~~~~~~~~~~~~~~~~~~~~ A frm




% ~~~~~~~~~~~~~~~~~~~~~~~~~~~~~~~~~~~~~~ V
\hSection[End]{Summary}
% ~~~~~~~~~~~~~~~~~~~~~~~~~~~~~~~~~~~~~~ A




% ~~~~~~~~~~~~~~~~~~~~~~~~~~~~~~~~~~~~~~ V frm
\begin{frame}[t,fragile] \frametitle{Cheat sheet}


% +++++++++++++++++++++++++++++++++++++++ V tbl
%: -Tbl. cheat sheet
\begin{table}
		\scriptsize
		\begin{tabular}{@{}L{0.30\textwidth}L{0.62\textwidth}@{}}
			\toprule
			\textbf{Topic} & \textbf{Remember} \\
			\midrule
			cost             & the stack trace makes exceptions slow; check first when bad input is common \\
			\hCode{finally}   & never \hCode{return} or \hCode{throw} from it \\
			resources        & try-with-resources keeps the first problem, \hCode{getSuppressed()} the others \\
			lambdas          & wrap checked exceptions, e.g. in \hCode{UncheckedIOException}, keeping the cause \\
			\hCode{Error}     & do not catch; fix the code or the configuration (\hCode{-Xss}, \hCode{-Xmx}) \\
			layers           & translate at the border, with the cause; handle once at the top \\
			\hCode{Optional}  & when ``no value'' is a normal answer \\
			test doubles     & inject the dependency; a stub or a spy is often one lambda \\
			\hCode{@TempDir}  & real files in a fresh folder for every test \\
			coverage         & finds code no test runs; not a quality measure \\
			mutation testing & how many small bugs do the tests notice? \\
			property-based   & many random inputs, one property, a fixed seed \\
			\bottomrule
		\end{tabular}
\end{table}%
% +++++++++++++++++++++++++++++++++++++++ A tbl
\end{frame}
% ~~~~~~~~~~~~~~~~~~~~~~~~~~~~~~~~~~~~~~ A frm




% ~~~~~~~~~~~~~~~~~~~~~~~~~~~~~~~~~~~~~~ V frm
\begin{frame}[t,fragile] \frametitle{Projects}
	\begin{enumerate}\small
		\item Measure \hCode{ExceptionCost} with 1\%, 10\% and 90\% bad words. Where does catch-and-parse become slower than check-first?
		\item Give \hCode{Report} a second dependency, a \hCode{Clock}, and print the date of the report. Test it with a fixed clock
		\item Add four mutants to \hCode{mutate.sh} that \hCode{StrongTest} does not kill, if you can; then add the tests that kill them
		\item Write a property-based test for your own \hCode{reverse(int[])}: reversing twice gives the original
		\item Design the exceptions of a three-layer grade program: which types, where translated, what the user sees
	\end{enumerate}
\end{frame}
% ~~~~~~~~~~~~~~~~~~~~~~~~~~~~~~~~~~~~~~ A frm




% ~~~~~~~~~~~~~~~~~~~~~~~~~~~~~~~~~~~~~~ V frm
\begin{frame}[t,fragile] \frametitle{References}
	\small The Java Tutorials, lesson Exceptions \cite{TheJavaTutorials}; JUnit User Guide \cite{homeJunit}; Big Java \cite{horstmann2016big}.

	\vspace{1em}
	\setbeamertemplate{bibliography item}[text]
	\bibliographystyle{IEEEtran}
	\tiny\bibliography{../../hblib/hblibDataStructures}
\end{frame}
% ~~~~~~~~~~~~~~~~~~~~~~~~~~~~~~~~~~~~~~ A frm




% ~~~~~~~~~~~~~~~~~~~~~~~~~~~~~~~~~~~~~~ V frm
\begin{frame}[t,fragile]
	\begin{center}
		\vspace{4cm}
		\hRemark{\Huge Questions?}
	\end{center}
\end{frame}
% ~~~~~~~~~~~~~~~~~~~~~~~~~~~~~~~~~~~~~~ A frm



\end{document}
